\section{Trả lời câu hỏi nghiên cứu}

Hai lớp kết quả được tách rõ. Phần lịch sử mô tả RTT/TCP/fan-in; phần xác nhận
mới dùng cùng cohort và block cân bằng để đánh giá RTT, sau đó dùng workload
HTTP để kiểm tra ảnh hưởng của khả năng phục vụ ứng dụng.

\subsection{RQ1: RTT khác nhau như thế nào?}

Trong môi trường thử nghiệm hiện tại, RTT mean qua Peering là 0,2078 ms và qua
TGW là 0,5630 ms. Peering thấp hơn 0,3552 ms. Đây là kết quả end-to-end của
cùng endpoint sau khi đổi route, không phải tuyên bố Peering luôn có latency
thấp hơn TGW.

\subsection{RQ2: TCP single-flow khác nhau như thế nào?}

Receiver goodput mean đạt 4,782804 Gbps qua Peering và 4,608265 Gbps qua TGW.
Theo định nghĩa comparison của independent review, Peering cao hơn 3,7875\%
trong các phiên đo hiện tại. Retransmission khác biệt lớn giữa hai phiên nhưng
chưa đủ bằng chứng để quy nguyên nhân cho riêng network service.

\subsection{RQ3: Tổng fan-in thay đổi như thế nào?}

Tổng goodput xấp xỉ đạt 11,611009 Gbps qua Peering và 11,590520 Gbps qua TGW,
chênh lệch 0,17646\% theo Peering làm mốc. Hai giá trị gần nhau trong dữ liệu
hiện có, nhưng chưa có equivalence test và cửa sổ source chưa trùng tuyệt đối.

\subsection{RQ4: Endpoint ảnh hưởng cách diễn giải ra sao?}

CPU busy B khoảng 59,18\% và 56,88\%; đồng thời bandwidth/PPS allowance
counter tại B tăng trong cả hai phiên fan-in. Vì ba source dùng chung receiver,
tổng goodput có thể bị giới hạn bởi EC2/ENA tại B. Fan-in vì vậy không cô lập
riêng khả năng chuyển tiếp của Peering hoặc TGW.

\subsection{RQ5: Trade-off performance và scalability là gì?}

Peering có RTT và single-flow goodput tốt hơn trong cấu hình đã đo, đồng thời
phù hợp khi số VPC nhỏ và cần kết nối trực tiếp. TGW đổi lại bằng topology
hub-and-spoke, định tuyến tập trung và số attachment tăng tuyến tính. Lựa chọn
phụ thuộc latency target, quy mô VPC, governance, cost và khả năng vận hành.

\subsection{RQ6: Network hay worker chi phối deadline?}

Ở 600 request/s, mọi điều kiện đạt 100\% đúng hạn trong ba block. Ở 800
request/s, cả Peering và TGW đều chỉ đạt khoảng 3\% đúng hạn với 8 slots, nhưng
đạt 100\% trong dữ liệu quan sát với 16 slots. Mean hiệu ứng 16 trừ 8 slots là
96,8417 điểm phần trăm. Queue audit cho thấy 8 slots tạo hàng trăm lời gọi chờ
trung bình, trong khi 16 slots gần như không có backlog. Vì vậy, trong workload
này, giới hạn đồng thời của worker là yếu tố chi phối deadline tại tải 800
request/s; chênh lệch route vẫn tồn tại nhưng nhỏ hơn hiệu ứng hàng đợi.

\section{Đánh giá trade-off kiến trúc}

Peering không bắc cầu. Full mesh bốn VPC cần sáu connection; khi thêm VPC, số
quan hệ và route cần quản trị tăng nhanh. Lợi thế là topology trực tiếp và
không có phí attachment theo giờ.

TGW thêm một lớp hub, VPC attachment, TGW route table và chi phí xử lý dữ liệu.
Đổi lại, network team có điểm kiểm soát tập trung, tách policy tốt hơn và không
cần tạo peering cho từng cặp. Đây là architectural advantage, không phải bằng
chứng TGW có throughput hoặc latency tốt hơn.

\section{Hạn chế}

\begin{itemize}
  \item Mỗi bộ dữ liệu chính chỉ có năm lượt trong một phiên.
  \item Đợt xác nhận dừng ở năm block RTT và ba block ứng dụng do ngân sách;
        thiết kế 16 block ban đầu chưa hoàn thành.
  \item Peering và TGW được đo nối tiếp, chưa counterbalance/randomize đầy đủ.
  \item Burst bandwidth, placement và nhiễu nền của EC2 có thể ảnh hưởng.
  \item Fan-in cộng mean từ các cửa sổ riêng, chưa căn common overlap window.
  \item ENA delta của B bao phủ cả phiên, chưa gắn cho từng lượt.
  \item Receiver B có dấu hiệu bandwidth/PPS allowance bottleneck.
  \item Phạm vi chỉ gồm một Region, một AZ và một loại/kích thước instance.
  \item Workload chỉ gồm idle RTT, TCP single-flow và TCP fan-in baseline.
  \item Workload HTTP dùng I/O mô phỏng 10 ms và hai worker; chưa đại diện cho
        database, cache hoặc microservice sản xuất.
  \item P95 tính trên request thành công; SLI deadline dùng toàn bộ request theo
        lịch. Request trong cùng phiên có thể phụ thuộc nhau.
  \item Trace Grafana là smoke lịch sử, không phải trace benchmark 800 rps;
        host metrics OpenTelemetry trên EC2 chưa được triển khai.
  \item Kết quả chưa xác định performance tối đa hoặc equivalence của dịch vụ.
\end{itemize}

\section{Hướng phát triển}

Các hạng mục sau là future work, chưa có kết quả thực nghiệm trong báo cáo:

\begin{itemize}
  \item nếu ngân sách cho phép, hoàn thành repeated balanced trials nhiều ngày;
  \item đo RTT idle và RTT under load, báo cáo P50/P95 và ECDF;
  \item phân tích TCP congestion window và RTT theo thời gian;
  \item tính fan-in trên common overlapping window;
  \item thử receiver lớn hơn để kiểm tra giả thuyết endpoint bottleneck;
  \item trace lấy mẫu có kiểm soát ở tải cao và bổ sung host metrics;
  \item xây dựng cost/performance analysis bằng bảng giá chốt theo ngày.
\end{itemize}

\section{Kết luận}

Trong môi trường bốn EC2 \texttt{c6i.large} cùng AZ, MTU 1500 và workload đã
đo, Peering cho thấy lợi thế về RTT idle và receiver goodput TCP một luồng.
TGW đem lại lợi thế về topology hub-and-spoke, định tuyến tập trung và khả năng
mở rộng quan hệ kết nối khi số VPC tăng.

Fan-in đạt tổng xấp xỉ gần nhau, nhưng network allowance counter tại B tăng ở
cả hai mode. Kết quả phù hợp với khả năng endpoint chi phối aggregate
throughput, không đủ để tuyên bố Peering và TGW có performance giống nhau.
Lựa chọn giải pháp cần xét đồng thời latency, scale, governance, cost và
workload; không nên dựa riêng vào một giá trị throughput.

Đợt workload ứng dụng bổ sung một kết luận quan trọng: network path có RTT thấp
không tự bảo đảm request đúng hạn. Khi 8 slots vận hành sát công suất danh
nghĩa, backlog làm P95 tăng lên hơn một giây qua cả hai route; tăng lên 16 slots
đưa P95 về khoảng 13 ms và toàn bộ request quan sát trở lại deadline. Vì vậy
quyết định kiến trúc phải xét đồng thời đường mạng, khả năng phục vụ, hàng đợi
và SLI của ứng dụng. Đây là kết luận có điều kiện trong cấu hình đã đo, không
phải tuyên bố Peering hoặc TGW thắng tổng quát.
